% ------------------------------------------------------------------------------
% Bab II - Landasan Teori
% ------------------------------------------------------------------------------
\chapter{LANDASAN TEORI}

% ------------------------------------------------------------------------------
% Tinjauan Pustaka
% ------------------------------------------------------------------------------
\section{Tinjauan Pustaka}

Tinjauan pustaka ini merangkum penelitian terdahulu yang berkaitan dengan reliabilitas \textit{LLM-as-a-Judge} sebagai acuan dan pembanding bagi Tugas Akhir ini. Tabel~\ref{tab:penelitian_terdahulu} menyajikan sepuluh penelitian terdahulu pada rentang tahun 2023--2026 yang relevan dengan tiga fokus utama penelitian, yaitu bias posisi jawaban (\textit{position bias}), sensitivitas terhadap variasi \textit{prompt} (\textit{prompt sensitivity}), serta evaluasi pada Bahasa Indonesia atau lintas bahasa. Setiap penelitian diuraikan berdasarkan pembahasan, metode, kelebihan dan kekurangan, serta perbedaannya dengan penelitian ini.

{\footnotesize
\linespread{1.0}\selectfont
\renewcommand{\arraystretch}{1.15}
\setlength{\tabcolsep}{4pt}
\setlength{\LTcapwidth}{\textwidth}
\setlength{\LTpre}{6pt}
\setlength{\LTpost}{6pt}
\begin{longtable}{|>{\centering\arraybackslash}p{0.55cm}|>{\raggedright\arraybackslash}p{2.05cm}|>{\raggedright\arraybackslash}p{2.55cm}|>{\raggedright\arraybackslash}p{2.10cm}|>{\raggedright\arraybackslash}p{2.45cm}|>{\raggedright\arraybackslash}p{2.40cm}|}
\caption{Ringkasan penelitian terdahulu.}\label{tab:penelitian_terdahulu}\\
\hline
\textbf{No.} &
\multicolumn{1}{>{\centering\arraybackslash}p{2.05cm}|}{\textbf{Judul dan Tahun}} &
\multicolumn{1}{>{\centering\arraybackslash}p{2.55cm}|}{\textbf{Pembahasan}} &
\multicolumn{1}{>{\centering\arraybackslash}p{2.10cm}|}{\textbf{Metode}} &
\multicolumn{1}{>{\centering\arraybackslash}p{2.45cm}|}{\textbf{Kelebihan dan Kekurangan}} &
\multicolumn{1}{>{\centering\arraybackslash}p{2.40cm}|}{\textbf{Perbedaan Penelitian}} \\
\hline
\endfirsthead

\multicolumn{6}{r}{\textit{(Lanjutan Tabel~\ref{tab:penelitian_terdahulu})}} \\
\hline
\textbf{No.} &
\multicolumn{1}{>{\centering\arraybackslash}p{2.05cm}|}{\textbf{Judul dan Tahun}} &
\multicolumn{1}{>{\centering\arraybackslash}p{2.55cm}|}{\textbf{Pembahasan}} &
\multicolumn{1}{>{\centering\arraybackslash}p{2.10cm}|}{\textbf{Metode}} &
\multicolumn{1}{>{\centering\arraybackslash}p{2.45cm}|}{\textbf{Kelebihan dan Kekurangan}} &
\multicolumn{1}{>{\centering\arraybackslash}p{2.40cm}|}{\textbf{Perbedaan Penelitian}} \\
\hline
\endhead

\hline
\multicolumn{6}{r}{\textit{(Bersambung ke halaman berikutnya)}} \\
\endfoot

\hline
\endlastfoot

1 & \textit{Judging LLM-as-a-Judge with MT-Bench and Chatbot Arena} (2023) \cite{zhengJudgingLLMasaJudgeMTBench2023} & Penelitian ini menguji penggunaan LLM sebagai penilai dialog terbuka. GPT-4 mencapai kesepakatan $>80\%$ dengan manusia pada pasangan penilaian yang tidak berakhir seri (\textit{non-tie}), tetapi konsistensi posisi \textit{prompt} standar hanya 23,8\%--65,0\% pada pasangan respons yang setara. & \textit{Pairwise comparison}, \textit{single-answer}, dan \textit{reference-guided grading} pada GPT-4, GPT-3.5, serta Claude-v1 dengan mitigasi \textit{swap-tie}. & \textbf{Kelebihan:} Divalidasi terhadap anotator ahli dan data \textit{crowdsourcing} berskala besar.\newline\textbf{Kekurangan:} Berbahasa Inggris, hanya memakai kecocokan mentah, dan aturan \textit{tie} menutupi besarnya \textit{flip rate}. & Menguji respons Bahasa Indonesia IndoPref pada \texttt{Qwen/\allowbreak Qwen3-\allowbreak 32B} dengan keluaran biner, Cohen's kappa, serta variasi \textit{prompt}. \\
\hline

2 & \textit{Split and Merge: Aligning Position Biases in LLM-based Evaluators} (2024) \cite{liSplitMergeAligning2024} & Penelitian ini mengkaji bias posisi pada 11.520 evaluasi pasangan MT-Bench. Konsistensi awal bervariasi (Claude2 28,28\% \textit{relation-based}; GPT-4 61,50\% \textit{likert-based} dengan $>78\%$ kasus tidak konsisten memilih posisi kedua), dan meningkat relatif 47,46\% dengan PORTIA. & Metode \textit{split-and-merge} (PORTIA) pada enam model \textit{judge} dengan format \textit{relation-based}, \textit{score-based}, dan \textit{likert-based}. & \textbf{Kelebihan:} Menunjukkan pengaruh format penilaian terhadap arah bias tanpa \textit{fine-tuning}.\newline\textbf{Kekurangan:} Jumlah token masukan bertambah panjang, gagal 17,13\% pada tugas kode, dan khusus bahasa Inggris. & Tidak memecah teks jawaban, melainkan menguji pengaruh posisi jawaban (A--B dan B--A) serta tiga variasi \textit{prompt} Bahasa Indonesia. \\
\hline

3 & \textit{Judging the Judges: A Systematic Study of Position Bias in LLM-as-a-Judge} (2025) \cite{shiJudgingJudgesSystematic} & Penelitian ini menganalisis \textit{position bias} pada $>150.000$ evaluasi MT-Bench dan DevBench. Stabilitas pengulangan umumnya $>0,85$ pada model \textit{judge} yang andal sehingga menunjukkan bias bersifat sistematis, sedangkan konsistensi posisi (0,57--0,82 di MT-Bench) turun saat kualitas kedua jawaban hampir setara. & Pertukaran urutan posisi \textit{pairwise} dan \textit{listwise} pada 15 \textit{judge} memakai \textit{repetition stability}, \textit{position consistency}, \textit{preference fairness}, dan regresi. & \textbf{Kelebihan:} Memisahkan bias sistematis dari variasi acak serta mengukur pengaruh selisih kualitas jawaban.\newline\textbf{Kekurangan:} Tidak menguji variasi \textit{prompt} dan terbatas pada bahasa Inggris. & Menggabungkan pengujian posisi jawaban dengan tiga variasi \textit{prompt} Bahasa Indonesia yang bermakna sama terhadap acuan manusia IndoPref. \\
\hline

4 & \textit{Diagnosing Bias and Instability in LLM Evaluation: A Scalable Pairwise Meta-Evaluator} (2025) \cite{anghelDiagnosingBiasInstability2025} & Penelitian ini menguji 3.600 perbandingan (1.800 pasangan bolak-balik A--B dan B--A). Sebanyak 48,4\% keputusan berbalik saat posisi jawaban ditukar meskipun tiga \textit{judge} bersepakat 100\%, dan penambahan instruksi anti-bias hanya memberi perbaikan terbatas. & Pengujian berpasangan dua arah, perubahan kata serta parafrase, dan penambahan aturan anti-bias pada LLaMA 3:8B, OpenHermes, serta Nous-Hermes 2. & \textbf{Kelebihan:} Menunjukkan bahwa kesepakatan 100\% antar-\textit{judge} tetap dapat menyimpan bias posisi yang sama.\newline\textbf{Kekurangan:} Hanya memakai 10 soal bahasa Inggris tanpa label acuan manusia. & Mengevaluasi 519 sampel unik dari 10 domain Bahasa Indonesia yang memiliki label \textit{reference human preference}. \\
\hline

5 & \textit{Systematic Evaluation of LLM-as-a-Judge in LLM Alignment Tasks: Explainable Metrics and Diverse Prompt Templates} (2025) \cite{weiSystematicEvaluationLLMasaJudge2025} & Penelitian ini menguji variasi \textit{template prompt} pada tugas peringkasan dan dialog. Akurasi dua urutan berada di bawah $0,7$ (turun hingga $<0,2$ pada \textit{template} yang kurang sesuai), dan pergantian \textit{template} dapat membalik arah \textit{position bias}. & Pertukaran posisi jawaban dengan koreksi \textit{flipping noise} pada tiga model OpenAI memakai 8 \textit{template} TL;DR dan 10 \textit{template} HH-RLHF. & \textbf{Kelebihan:} Menunjukkan hubungan antara pilihan \textit{template}, arah bias posisi, dan akurasi.\newline\textbf{Kekurangan:} Hanya menguji model tertutup OpenAI pada data berbahasa Inggris. & Menguji model terbuka (\textit{open-weight}) \texttt{Qwen/\allowbreak Qwen3-\allowbreak 32B} pada Bahasa Indonesia dengan tiga variasi \textit{prompt} yang bermakna sama. \\
\hline

6 & \textit{JudgeSense: A Benchmark for Prompt Sensitivity in LLM-as-a-Judge Systems} (2026) \cite{bellibatluJudgeSenseBenchmarkPrompt2026} & Penelitian ini mengukur sensitivitas model terhadap parafrase instruksi yang bermakna sama. Perubahan kalimat menurunkan konsistensi di empat tugas, skala ordinal 1--5 jauh kurang stabil dibandingkan biner, dan ukuran parameter dalam satu keluarga model tidak menjamin stabilitas. & Pengujian 880 sampel berlabel manusia pada 25 \textit{judge} (termasuk Qwen3) melalui selisih kesepakatan parafrase dan pengulangan ($\Delta\text{JSS}$). & \textbf{Kelebihan:} Memisahkan pengaruh perubahan kalimat instruksi dari pengaruh pengulangan inferensi.\newline\textbf{Kekurangan:} Terbatas pada bahasa Inggris dan hanya menilai respons tunggal (\textit{pointwise}) tanpa menguji pertukaran posisi jawaban. & Menguji perbandingan berpasangan (\textit{pairwise}) Bahasa Indonesia IndoPref dengan tiga variasi \textit{prompt} yang bermakna sama (P1, P2, P3) pada urutan posisi A--B dan B--A. \\
\hline

7 & \textit{The Coin Flip Judge? Reliability and Bias in LLM-as-a-Judge Evaluation} (2026) \cite{yagubyanCoinFlipJudge2026} & Penelitian ini menguji kestabilan keputusan pada pasangan jawaban yang setara. Hasil pengulangan berbalik rata-rata 13,6\%, GPT-4o-mini memilih posisi pertama pada 72\% soal, muncul \textit{pairwise--pointwise gap}, dan \textit{template} yang bermakna sama mengubah pemenang mayoritas pada 25\% kasus. & Pengulangan 50 \textit{trial} \textit{pairwise} dan \textit{pointwise} pada GPT-4o-mini serta GPT-4.1-mini di 29 soal, disertai pengujian $T=0$ dan dua \textit{template}. & \textbf{Kelebihan:} Menunjukkan ketidakkonsistenan pada jawaban yang setara yang tetap muncul pada $T=0$.\newline\textbf{Kekurangan:} Hanya menguji dua model OpenAI pada 29 soal bahasa Inggris tanpa label manusia. & Mengevaluasi 519 sampel Bahasa Indonesia berlabel preferensi manusia pada model \textit{open-weight} dengan tiga variasi \textit{prompt} dan dua posisi jawaban. \\
\hline

8 & \textit{Reliability without Validity: A Systematic, Large-Scale Evaluation of LLM-as-a-Judge Models Across Agreement, Consistency, and Bias} (2026) \cite{normanReliabilityValiditySystematic2026} & Penelitian ini mengevaluasi 21 \textit{judge} ($\sim 541.000$ penilaian) dan menemukan \textit{kappa deflation} sebesar 33,8--41,3 poin persentase di MT-Bench serta \textit{consistency--bias paradox} pada model dengan konsistensi pengulangan Krippendorff's $\alpha > 0,95$ yang tetap memiliki bias posisi tinggi. & Pengujian \textit{agreement}, \textit{consistency}, dan \textit{bias audit} ($T=0$) di MT-Bench, JudgeBench, serta RewardBench memakai \textit{exact match}, Cohen's kappa, dan \textit{flip rate}. & \textbf{Kelebihan:} Menunjukkan pentingnya penggunaan Cohen's kappa bersama pengujian posisi jawaban.\newline\textbf{Kekurangan:} Hanya memakai satu \textit{template} bahasa Inggris tanpa variasi \textit{prompt}. & Menerapkan \textit{accuracy}, Cohen's kappa, dan pengujian posisi jawaban pada Bahasa Indonesia sekaligus menguji tiga variasi \textit{prompt}. \\
\hline

9 & \textit{How Reliable is Multilingual LLM-as-a-Judge?} (2025) \cite{fuHowReliableMultilingual2025} & Penelitian ini menguji konsistensi penilaian lintas 25 bahasa paralel (termasuk 142 sampel Bahasa Indonesia pada WikiLingua). Rata-rata konsistensi masih rendah (Fleiss' kappa $\approx 0,3$), format biner lebih stabil daripada skala 1--5, dan ukuran model tidak menjamin konsistensi. & Evaluasi \textit{pointwise} (biner \textit{Yes/No} dan skala 1--5) pada lima \textit{judge} di lima dataset paralel memakai Fleiss' kappa, Cohen's kappa, dan \textit{ensemble}. & \textbf{Kelebihan:} Mengukur ketidakkonsistenan lintas bahasa dan menunjukkan kestabilan format biner.\newline\textbf{Kekurangan:} Hanya menguji skema \textit{pointwise} dengan \textit{prompt} bahasa Inggris tanpa pertukaran posisi. & Berfokus pada perbandingan \textit{pairwise} dengan teks jawaban, \textit{prompt}, dan preferensi manusia berbahasa Indonesia. \\
\hline

10 & \textit{IndoPref: A Multi-Domain Pairwise Preference Dataset for Indonesian} (2025) \cite{wiyonoIndoPrefMultiDomainPairwise2025} & Penelitian ini membangun dataset preferensi berpasangan Bahasa Indonesia (522 \textit{prompt}, 4.099 pasangan lintas 10 domain, Krippendorff's $\alpha = 0,891$) serta menguji 15 model. Akurasi tertinggi mencapai 74,34\%, dan terjemahan \textit{prompt} ke Bahasa Indonesia menurunkan akurasi GPT-4.1 sebesar 6,39 poin persentase. & Penilaian anotator berbahasa Indonesia dikonversi menjadi preferensi berpasangan, lalu diuji akurasinya pada satu urutan menggunakan 12 model \textit{open-weight} dan 3 model tertutup. & \textbf{Kelebihan:} Menyediakan dataset \textit{human preference} asli Bahasa Indonesia lintas 10 domain.\newline\textbf{Kekurangan:} Hanya mengukur akurasi satu urutan tanpa pertukaran posisi A--B dan B--A, tanpa Cohen's kappa, serta tanpa variasi \textit{prompt}. & Menggunakan IndoPref untuk menguji reliabilitas \texttt{Qwen/\allowbreak Qwen3-\allowbreak 32B} pada pertukaran posisi jawaban (A--B dan B--A), tiga variasi \textit{prompt} yang bermakna sama, dan Cohen's kappa. \\
\end{longtable}
}

Perbandingan penelitian pada Tabel~\ref{tab:penelitian_terdahulu} beserta literatur terkait menunjukkan bahwa evaluasi \textit{LLM-as-a-Judge} tidak cukup hanya dilihat dari persentase kecocokan rata-rata, melainkan perlu mengukur kesepakatan terhadap penilaian manusia sekaligus konsistensi keputusan model itu sendiri. Penelitian awal melaporkan bahwa model penilai dapat mencapai kecocokan di atas 80\% dengan manusia pada dialog terbuka \cite{zhengJudgingLLMasaJudgeMTBench2023}, tetapi pengujian lintas 20 tugas pemrosesan bahasa alami menunjukkan tingkat kesepakatan yang sangat bervariasi antartugas \cite{bavarescoLLMsInsteadHuman2025}, bahkan mendekati tebakan acak pada tugas verifikasi objektif ketika diuji dalam dua urutan posisi jawaban \cite{tanJudgeBenchBenchmarkEvaluating2025}. Dari sisi metrik evaluasi, persentase kecocokan mentah belum cukup menggambarkan kualitas penilai karena model dengan kecocokan rata-rata di atas 90\% masih dapat berselisih lebih dari 10 poin persentase terhadap skor manusia \cite{thakurJudgingJudgesEvaluating2025}. Selain itu, nilai kecocokan mentah terbukti lebih tinggi 33,8--41,3 poin persentase dibandingkan Cohen's kappa pada MT-Bench (\textit{kappa deflation}) karena belum memperhitungkan peluang kebetulan \cite{normanReliabilityValiditySystematic2026}, serta mudah berubah akibat perbedaan cara menangani kasus ketika model tidak memberikan pilihan jawaban, keluaran \textit{invalid}, dan ketimpangan proporsi label \cite{raoAgreementMetricsLLMasJudge2026}. Temuan bahwa tiga model \textit{judge} yang bersepakat 100\% tetap membalik 48,4\% keputusannya saat posisi jawaban ditukar \cite{anghelDiagnosingBiasInstability2025}, serta tingginya konsistensi pengulangan (Krippendorff's $\alpha = 0,992$ di MT-Bench) pada model yang memiliki bias posisi besar seperti Qwen 3 8B \cite{normanReliabilityValiditySystematic2026}, turut menunjukkan adanya \textit{consistency--bias paradox}, yaitu keluaran yang deterministik belum tentu bebas dari bias sistematis.

Hasil penelitian terdahulu juga menunjukkan bahwa arah \textit{position bias} dan \textit{prompt sensitivity} berbeda-beda antarmodel maupun antarpengaturan uji. Beberapa penelitian menemukan kecenderungan kuat untuk memilih jawaban pada posisi pertama (\textit{primacy bias}), seperti 75,0\% pada Claude-v1 \cite{zhengJudgingLLMasaJudgeMTBench2023} dan 72\% kemenangan mayoritas posisi pertama pada GPT-4o-mini \cite{yagubyanCoinFlipJudge2026}, sedangkan penelitian lain melaporkan kecenderungan memilih jawaban pada posisi kedua (\textit{recency bias}), seperti lebih dari 78\% kasus tidak konsisten pada GPT-4 dengan format Likert \cite{liSplitMergeAligning2024} serta perbedaan tanda nilai \textit{preference fairness} antarmodel \cite{shiJudgingJudgesSystematic}. Ketidakkonsistenan akibat posisi ini semakin besar ketika kualitas kedua jawaban hampir setara \cite{shiJudgingJudgesSystematic}\cite{yagubyanCoinFlipJudge2026}. Selain itu, bias posisi berkaitan erat dengan penyusunan \textit{prompt} karena pergantian \textit{template} pada model yang sama dapat menggeser akurasi sekaligus membalik arah \textit{position bias} \cite{weiSystematicEvaluationLLMasaJudge2025}. Bahkan ketika instruksi ditulis ulang dengan makna yang sama (\textit{semantically equivalent}), konsistensi keputusan tetap turun melebihi pengaruh pengulangan inferensi (\textit{decoding noise}) \cite{bellibatluJudgeSenseBenchmarkPrompt2026} dan mengubah pemenang mayoritas pada 25\% kasus \cite{yagubyanCoinFlipJudge2026}, sementara ukuran parameter yang lebih besar dalam satu keluarga model tidak menjamin kestabilan keputusan terhadap perubahan kalimat \textit{prompt} \cite{bellibatluJudgeSenseBenchmarkPrompt2026}.

Meskipun berbagai sumber ketidakstabilan tersebut telah banyak diteliti, sebagian besar pengujian masih berfokus pada bahasa Inggris. Hanya 33 dari 650 publikasi \textit{LLM-as-a-Judge} (5,1\%) di ACL Anthology yang meneliti konteks multibahasa atau bahasa dengan keterbatasan data (\textit{low-resource}), dan banyak di antaranya belum disertai validasi manusia yang memadai pada bahasa yang diuji \cite{dogruozChallengesRecommendationsLLMasaJudge2026}. Padahal, pengujian lintas bahasa menunjukkan konsistensi penilaian yang rendah dengan rata-rata Fleiss' kappa sekitar 0,3 \cite{fuHowReliableMultilingual2025}, skor penilaian yang cenderung mengumpul di nilai tengah (\textit{score convergence}) pada bahasa \textit{low-resource} \cite{sonMMEvalMultilingualMetaEvaluation2025}, perbedaan akurasi antarbahasa serta bias terhadap bahasa jawaban tertentu \cite{zhouFairnessFluencyInvestigation2026}, dan meningkatnya ketidakkonsistenan posisi ketika bahasa pengujian bergeser dari bahasa Inggris \cite{yinDoesJudgePrefer2026}. Penelitian multibahasa yang menyertakan Bahasa Indonesia umumnya baru menggunakan data ringkasan WikiLingua dengan skema \textit{pointwise} \cite{fuHowReliableMultilingual2025}, data hasil terjemahan mesin \cite{pombalMPrometheusSuiteOpen2025}, atau pasangan opsi soal pilihan ganda terjemahan MMMLU \cite{zhouFairnessFluencyInvestigation2026}. Sementara itu, pengujian pada dataset preferensi berpasangan Bahasa Indonesia IndoPref \cite{wiyonoIndoPrefMultiDomainPairwise2025} baru mengukur akurasi pada satu urutan jawaban dan satu terjemahan \textit{prompt} tanpa menguji pertukaran posisi jawaban maupun variasi \textit{prompt}. Oleh karena itu, meskipun keberadaan \textit{position bias} dan \textit{prompt sensitivity} pada \textit{LLM-as-a-Judge} telah banyak dibuktikan pada bahasa Inggris, penelitian khusus Bahasa Indonesia yang mengacu pada \textit{human preference} serta menguji pengaruh posisi jawaban dan variasi \textit{prompt} secara terpisah masih terbatas.

% ------------------------------------------------------------------------------
% Landasan Teori
% ------------------------------------------------------------------------------
\section{Landasan Teori}

Bagian ini menguraikan konsep dan dasar teori yang digunakan dalam pelaksanaan penelitian serta analisis pada bab-bab berikutnya. Pembahasan mencakup konsep \textit{large language model} (LLM), metode \textit{LLM-as-a-Judge}, ukuran reliabilitas model penilai, \textit{position bias}, \textit{prompt sensitivity}, evaluasi LLM pada Bahasa Indonesia, serta sintesis teori dan celah penelitian.

\subsection{Large Language Models}

\textit{Large language model} (LLM) adalah model bahasa generatif berbasis jaringan saraf tiruan berskala besar yang memodelkan distribusi probabilitas kondisional atas barisan token secara autoregresif \cite{guSurveyLLMasaJudge2025}. Melalui tahap prapelatihan pada kumpulan data teks berukuran besar, model mempelajari pola bahasa dan pengetahuan lintas domain untuk memprediksi token berikutnya berdasarkan teks masukan. Agar dapat mengikuti perintah pengguna dengan tepat, model dilatih lebih lanjut melalui penyelarasan instruksi (\textit{instruction tuning}) dan preferensi manusia sehingga menghasilkan model \textit{instruction-tuned} yang mampu memahami instruksi peran, kriteria penilaian, serta aturan format keluaran secara langsung pada saat inferensi \cite{guSurveyLLMasaJudge2025}\cite{liLLMsasJudgesComprehensiveSurvey2024}.

Kemampuan tersebut memungkinkan pelaksanaan evaluasi berbasis inferensi (\textit{inference-based evaluation}) dengan menggunakan bobot parameter model secara langsung tanpa pelatihan ulang atau \textit{fine-tuning} \cite{guSurveyLLMasaJudge2025}\cite{liLLMsasJudgesComprehensiveSurvey2024}. Pada pendekatan ini, keputusan model ditentukan oleh interaksi antara bobot hasil pelatihan, susunan isi \textit{prompt} masukan, serta pengaturan dekoding seperti \textit{temperature} dan batas maksimum token keluaran \cite{bellibatluJudgeSenseBenchmarkPrompt2026}. Salah satu model terbuka (\textit{open-weight}) yang mendukung pendekatan ini adalah \texttt{Qwen/Qwen3-32B}, yaitu model \textit{dense} terbesar pada keluarga Qwen3 berukuran 32 miliar parameter yang dilatih pada 36 triliun token lintas 119 bahasa dan dialek termasuk Bahasa Indonesia \cite{yangQwen3TechnicalReport2025}. Model ini menggabungkan mode penalaran (\textit{thinking}) dan mode respons langsung (\textit{non-thinking}) dalam satu model yang sama, sehingga pengaktifan mode \textit{non-thinking} memungkinkan model menghasilkan label keputusan secara langsung tanpa tahapan penalaran internal \cite{yangQwen3TechnicalReport2025}. Pada pengujian \textit{benchmark} Bahasa Indonesia dalam mode \textit{non-thinking} yang mencakup INCLUDE, MMMLU, MT-AIME24, dan PolyMath, \texttt{Qwen/Qwen3-32B} memperoleh rata-rata skor 51,6 (dengan skor 81,8 pada INCLUDE dan 77,2 pada MMMLU), yang merupakan capaian tertinggi di antara seluruh model \textit{dense} yang diuji serta melampaui model pembanding seperti \texttt{Gemma-3-27b-IT} (47,0), \texttt{GPT-4o-2024-1120} (46,4), dan \texttt{Qwen2.5-72B-Instruct} (45,0) \cite{yangQwen3TechnicalReport2025}.

\subsection{LLM-as-a-Judge}

\textit{LLM-as-a-Judge} adalah metode evaluasi otomatis yang menggunakan LLM untuk menilai kualitas keluaran teks berdasarkan instruksi, kriteria, atau preferensi yang tertulis di dalam \textit{prompt} \cite{zhengJudgingLLMasaJudgeMTBench2023}\cite{guSurveyLLMasaJudge2025}. Mengadaptasi formalisasi evaluasi berbasis LLM \cite{guSurveyLLMasaJudge2025}\cite{liLLMsasJudgesComprehensiveSurvey2024}, apabila $x$ menyatakan instruksi tugas awal, $\mathcal{R}$ menyatakan pasangan jawaban yang dinilai, dan $\mathcal{C}$ menyatakan instruksi \textit{prompt} evaluasi, maka model memetakan seluruh masukan tersebut melalui distribusi probabilitas kondisional $\mathcal{P}_{\text{LLM}}$ dan fungsi ekstraksi keluaran $g(\cdot)$ menjadi label keputusan $\hat{y} = g(\mathcal{P}_{\text{LLM}}(x, \mathcal{R}, \mathcal{C}))$.

Metode penilaian ini secara umum terbagi menjadi skema penilaian tunggal (\textit{pointwise}) dan perbandingan berpasangan (\textit{pairwise comparison}) \cite{zhengJudgingLLMasaJudgeMTBench2023}\cite{liLLMsasJudgesComprehensiveSurvey2024}. Skema \textit{pointwise} menilai setiap respons secara terpisah dalam bentuk skor angka, tetapi rentan terhadap ketidakkonsistenan skala nilai \cite{liSplitMergeAligning2024}, menyempitnya rentang skor pada bahasa \textit{low-resource} \cite{sonMMEvalMultilingualMetaEvaluation2025}, serta variasi nilai pada soal yang sama sehingga menyulitkan pembedaan dua jawaban yang kualitasnya hampir setara \cite{zhengJudgingLLMasaJudgeMTBench2023}\cite{liLLMsasJudgesComprehensiveSurvey2024}\cite{yagubyanCoinFlipJudge2026}. Sebaliknya, \textit{pairwise comparison} membandingkan dua jawaban secara langsung untuk instruksi yang sama dan terbukti lebih selaras dengan preferensi manusia ketika perbedaan kualitas kedua jawaban cukup tipis \cite{zhengJudgingLLMasaJudgeMTBench2023}\cite{liLLMsasJudgesComprehensiveSurvey2024}. Pembatasan keluaran model menjadi pilihan biner (\texttt{A} atau \texttt{B}) juga menghindari ketidakkonsistenan penentuan batas nilai antarkategori yang terbukti jauh lebih besar pada penilaian berskala ordinal 1--5 \cite{bellibatluJudgeSenseBenchmarkPrompt2026}\cite{fuHowReliableMultilingual2025}. Namun, penyajian dua jawaban secara berurutan di dalam satu \textit{prompt} membuat skema berpasangan rentan terhadap bias posisi \cite{guSurveyLLMasaJudge2025} serta dapat memaksakan penentuan pemenang meskipun selisih skor kedua jawaban pada penilaian tunggal sangat kecil (\textit{pairwise--pointwise gap}) \cite{yagubyanCoinFlipJudge2026}.

\subsection{Reliabilitas LLM Judge}

Reliabilitas \textit{LLM-as-a-Judge} mencakup dua aspek utama, yaitu kesepakatan keputusan model dengan preferensi manusia (\textit{human agreement}) dan konsistensi keputusan model ketika cara penyajiannya diubah tanpa mengubah isi jawaban yang dinilai \cite{guSurveyLLMasaJudge2025}\cite{normanReliabilityValiditySystematic2026}. Kedua aspek ini perlu diukur bersama karena model yang memberikan hasil tetap saat diuji berulang pada masukan yang sama masih dapat menyimpan bias sistematis yang menyimpang dari preferensi manusia \cite{anghelDiagnosingBiasInstability2025}\cite{normanReliabilityValiditySystematic2026}, serta menggeser peringkat kualitas model apabila hasil penilaiannya digabungkan tanpa memperhitungkan keandalan penilai \cite{xuJudgeAwareRankingFramework2026}. Oleh karena itu, reliabilitas dalam penelitian ini diukur melalui kesepakatan terhadap \textit{human preference}, konsistensi terhadap posisi jawaban, dan konsistensi terhadap variasi \textit{prompt}.

Tingkat kesepakatan terhadap \textit{reference human preference} pertama-tama diukur menggunakan proporsi kecocokan keputusan valid (\textit{accuracy}) pada Persamaan~(\ref{eq:accuracy}):
\begin{equation}
  \text{Accuracy} = \frac{1}{N_{\text{valid}}} \sum_{i=1}^{N_{\text{valid}}} \mathbb{I}\left(\hat{y}_i = y_i\right)
  \label{eq:accuracy}
\end{equation}
dengan $N_{\text{valid}}$ adalah jumlah keputusan biner yang valid (\texttt{A} atau \texttt{B}), $\hat{y}_i \in \{\texttt{A}, \texttt{B}\}$ adalah keputusan \textit{judge}, $y_i \in \{\texttt{A}, \texttt{B}\}$ adalah label acuan manusia, dan $\mathbb{I}(\cdot)$ adalah fungsi indikator yang bernilai $1$ jika kondisi di dalamnya benar serta $0$ jika salah. Karena \textit{accuracy} belum memperhitungkan kesepakatan yang terjadi secara kebetulan \cite{raoAgreementMetricsLLMasJudge2026}\cite{thakurJudgingJudgesEvaluating2025} sehingga nilainya dapat jauh lebih tinggi daripada kesepakatan terkoreksi peluang (\textit{kappa deflation}) \cite{normanReliabilityValiditySystematic2026}, pengukurannya dilengkapi dengan Cohen's kappa ($\kappa$) \cite{normanReliabilityValiditySystematic2026}\cite{raoAgreementMetricsLLMasJudge2026} pada Persamaan~(\ref{eq:cohen_kappa}):
\begin{equation}
  \kappa = \frac{P_o - P_e}{1 - P_e}
  \label{eq:cohen_kappa}
\end{equation}
dengan $P_o = \text{Accuracy}$ dan $P_e = P(y = \texttt{A})P(\hat{y} = \texttt{A}) + P(y = \texttt{B})P(\hat{y} = \texttt{B})$ sebagai peluang kesepakatan kebetulan berdasarkan proporsi pilihan label manusia dan \textit{judge}. Mengingat nilai \textit{accuracy} maupun $\kappa$ sangat dipengaruhi oleh cara penanganan keluaran yang gagal (\textit{invalid}), keluaran model yang tidak menghasilkan label biner \texttt{A} atau \texttt{B} tidak dimasukkan ke salah satu kelas, melainkan dipisahkan dan dilaporkan persentasenya melalui \textit{invalid rate} \cite{raoAgreementMetricsLLMasJudge2026}.

\subsection{Position Bias}

\textit{Position bias} adalah kecenderungan sistematis model penilai untuk memilih salah satu jawaban berdasarkan urutan penempatannya di dalam \textit{prompt}, terlepas dari isi respons yang dinilai \cite{zhengJudgingLLMasaJudgeMTBench2023}\cite{shiJudgingJudgesSystematic}. Bias ini terjadi ketika pertukaran posisi jawaban mengubah keputusan akhir ($\text{LLM}(x, R_{\text{A}}, R_{\text{B}}) \neq \text{LLM}(x, R_{\text{B}}, R_{\text{A}})$) \cite{liSplitMergeAligning2024}, baik berupa kecenderungan memilih jawaban di posisi pertama (\textit{primacy bias}) maupun di posisi kedua (\textit{recency bias}) \cite{shiJudgingJudgesSystematic}, dan semakin kuat ketika kualitas kedua jawaban hampir setara \cite{liLLMsasJudgesComprehensiveSurvey2024}\cite{shiJudgingJudgesSystematic}.

Pengujian \textit{position bias} dilakukan melalui pertukaran posisi jawaban antara kondisi A--B (Respons A ditempatkan pada posisi pertama, Respons B pada posisi kedua) dan kondisi B--A (Respons B ditempatkan pada posisi pertama, Respons A pada posisi kedua) dengan tetap mempertahankan \texttt{A} serta \texttt{B} sebagai identitas jawaban asli \cite{zhengJudgingLLMasaJudgeMTBench2023}\cite{shiJudgingJudgesSystematic}\cite{anghelDiagnosingBiasInstability2025}. Tanpa mengubah keputusan yang berlawanan menjadi hasil seri (\textit{tie}), pengaruh posisi jawaban diukur menggunakan \textit{position consistency}, \textit{preference flip rate}, dan \textit{first-position preference rate} pada Persamaan~(\ref{eq:position_consistency})--(\ref{eq:first_pos_rate}) yang diadaptasi dari metrik evaluasi bias posisi \cite{shiJudgingJudgesSystematic}\cite{normanReliabilityValiditySystematic2026}\cite{yagubyanCoinFlipJudge2026}:
\begin{equation}
  \text{Position Consistency} = \frac{1}{N_{\text{pos}}} \sum_{i=1}^{N_{\text{pos}}} \mathbb{I}\left(\hat{y}_{i,\text{AB}} = \hat{y}_{i,\text{BA}}\right)
  \label{eq:position_consistency}
\end{equation}
\begin{equation}
  \text{Preference Flip Rate} = \frac{1}{N_{\text{pos}}} \sum_{i=1}^{N_{\text{pos}}} \mathbb{I}\left(\hat{y}_{i,\text{AB}} \neq \hat{y}_{i,\text{BA}}\right) = 1 - \text{Position Consistency}
  \label{eq:preference_flip_rate}
\end{equation}
\begin{equation}
  \text{First-Position Preference Rate} = \frac{1}{2N_{\text{pos}}} \sum_{i=1}^{N_{\text{pos}}} \Bigl[\mathbb{I}\left(\hat{y}_{i,\text{AB}} = \texttt{A}\right) + \mathbb{I}\left(\hat{y}_{i,\text{BA}} = \texttt{B}\right)\Bigr]
  \label{eq:first_pos_rate}
\end{equation}
dengan $N_{\text{pos}}$ menyatakan jumlah pasangan sampel yang memiliki keputusan valid pada kedua urutan, serta $\hat{y}_{i,\text{AB}}, \hat{y}_{i,\text{BA}} \in \{\texttt{A}, \texttt{B}\}$ menyatakan identitas jawaban asli yang dipilih. Pada pengujian dua urutan yang seimbang, \textit{first-position preference rate} digunakan untuk membedakan perubahan keputusan yang bersifat acak dari kecenderungan sistematis untuk memilih jawaban di posisi pertama (nilai $>0,50$) atau posisi kedua (nilai $<0,50$).

\subsection{Prompt Sensitivity}

\textit{Prompt sensitivity} adalah ketidakstabilan keputusan \textit{LLM-as-a-Judge} akibat perubahan kalimat instruksi, rincian kriteria penilaian, atau aturan tambahan di dalam \textit{prompt} meskipun tujuan evaluasi dan pasangan jawaban yang dinilai tetap sama \cite{guSurveyLLMasaJudge2025}\cite{weiSystematicEvaluationLLMasaJudge2025}\cite{bellibatluJudgeSenseBenchmarkPrompt2026}. Karena proses dekoding autoregresif bergantung pada susunan kata dan token di dalam \textit{prompt}, perbedaan struktur \textit{template} maupun variasi kalimat yang bermakna sama (\textit{semantically equivalent}) tetap dapat menggeser fokus penilaian model serta mengubah arah bias posisinya \cite{weiSystematicEvaluationLLMasaJudge2025}\cite{bellibatluJudgeSenseBenchmarkPrompt2026}.

Penyusunan \textit{prompt} evaluasi terstruktur dengan komponen peran, konteks, tugas penilaian, dan format keluaran umumnya mencakup tiga pendekatan, yaitu perbandingan menyeluruh secara langsung (\textit{baseline}) \cite{zhengJudgingLLMasaJudgeMTBench2023}\cite{guSurveyLLMasaJudge2025}, perincian rubrik multi-kriteria seperti akurasi, relevansi, dan kejelasan untuk memperjelas dasar penilaian \cite{guSurveyLLMasaJudge2025}\cite{liLLMsasJudgesComprehensiveSurvey2024}, serta penambahan aturan eksplisit anti-bias agar model mengabaikan urutan jawaban \cite{zhengJudgingLLMasaJudgeMTBench2023}\cite{anghelDiagnosingBiasInstability2025} dan panjang respons \cite{zhengJudgingLLMasaJudgeMTBench2023}. Penalaran bertahap (\textit{Chain-of-Thought}) atau penulisan penjelasan sebelum keputusan tidak digunakan pada pengujian utama agar keluaran model tetap berupa label biner, mengingat permintaan penjelasan bersamaan dengan keputusan terbukti dapat menurunkan konsistensi posisi \cite{guSurveyLLMasaJudge2025} serta penambahan instruksi penalaran bertahap dapat mengubah keputusan mayoritas model \cite{yagubyanCoinFlipJudge2026}. Pada posisi jawaban yang sama, konsistensi keputusan lintas tiga variasi \textit{prompt} (P1, P2, dan P3) diukur menggunakan \textit{prompt consistency} serta \textit{prompt flip rate} antarpasangan \textit{prompt} pada Persamaan~(\ref{eq:prompt_consistency}) dan Persamaan~(\ref{eq:prompt_flip_rate}) yang diadaptasi dari pengujian sensitivitas \textit{prompt} \cite{bellibatluJudgeSenseBenchmarkPrompt2026}\cite{yagubyanCoinFlipJudge2026}:
\begin{equation}
  \text{Prompt Consistency} = \frac{1}{N_{\text{prompt}}} \sum_{i=1}^{N_{\text{prompt}}} \mathbb{I}\left(\hat{y}_{i,\text{P1}} = \hat{y}_{i,\text{P2}} = \hat{y}_{i,\text{P3}}\right)
  \label{eq:prompt_consistency}
\end{equation}
\begin{equation}
  \text{Prompt Flip Rate}_{(p,q)} = \frac{1}{N_{(p,q)}} \sum_{i=1}^{N_{(p,q)}} \mathbb{I}\left(\hat{y}_{i,p} \neq \hat{y}_{i,q}\right)
  \label{eq:prompt_flip_rate}
\end{equation}
dengan $N_{\text{prompt}}$ menyatakan jumlah sampel yang menghasilkan keputusan valid pada ketiga variasi \textit{prompt} dan $N_{(p,q)}$ menyatakan jumlah sampel valid pada pasangan $(p,q) \in \{(\text{P1},\text{P2}), (\text{P1},\text{P3}), (\text{P2},\text{P3})\}$. Perhitungan \textit{flip rate} untuk setiap pasangan \textit{prompt} ini digunakan untuk melihat variasi \textit{prompt} mana yang paling banyak mengubah keputusan model.

\subsection{Evaluasi LLM Bahasa Indonesia}

Evaluasi otomatis pada teks Bahasa Indonesia masih menghadapi kendala karena terbatasnya porsi bahasa non-Inggris dalam data pelatihan LLM. Kondisi ini berdampak pada menyempitnya selisih skor penilaian pada bahasa dengan keterbatasan data (\textit{low-resource}) \cite{sonMMEvalMultilingualMetaEvaluation2025}, rendahnya konsistensi keputusan lintas bahasa \cite{fuHowReliableMultilingual2025}, serta munculnya perbedaan performa antarbahasa dan bias terhadap bahasa jawaban tertentu \cite{zhouFairnessFluencyInvestigation2026}. Dataset evaluasi yang diterjemahkan dari bahasa Inggris juga sering kali kehilangan kealamian kalimat serta konteks budaya lokal \cite{wiyonoIndoPrefMultiDomainPairwise2025}, sedangkan penggunaan data terjemahan mesin dalam pelatihan penilai multibahasa terbukti menurunkan ketepatan model pada data pengujian yang ditulis langsung dalam bahasa sasaran \cite{pombalMPrometheusSuiteOpen2025}. Di samping itu, penerjemahan instruksi \textit{prompt} evaluasi ke Bahasa Indonesia terbukti dapat menurunkan akurasi penilai \cite{wiyonoIndoPrefMultiDomainPairwise2025}, sementara perubahan bahasa penyajian dari bahasa Inggris memicu perubahan keputusan sebesar 10,7\%--14,4\% sekaligus meningkatkan ketidakkonsistenan posisi jawaban dari 12,1\% menjadi 14,9\%--18,1\% \cite{yinDoesJudgePrefer2026}.

Untuk menyediakan acuan evaluasi Bahasa Indonesia yang disusun dan dinilai langsung oleh anotator berbahasa Indonesia, dataset IndoPref memuat 522 \textit{prompt} lintas 10 domain (\textit{Analysis}, \textit{Brainstorming}, \textit{Coding}, \textit{Creative Writing}, \textit{Logic}, \textit{Math}, \textit{Open Question}, \textit{Safety}, \textit{Summarization}, dan \textit{Translation}) beserta 4.099 pasangan preferensi atas respons lima LLM \textit{instruction-tuned} dengan tingkat kesepakatan antar-anotator Krippendorff's $\alpha = 0,891$ \cite{wiyonoIndoPrefMultiDomainPairwise2025}. Dalam penelitian ini, label pada IndoPref digunakan sebagai \textit{reference human preference} sesuai prosedur anotasinya, bukan sebagai penilaian mutlak yang mewakili seluruh pengguna Bahasa Indonesia.

\subsection{Sintesis dan Research Gap}

Berdasarkan landasan teori yang telah diuraikan, ketika model \textit{instruction-tuned} digunakan sebagai \textit{pairwise LLM-as-a-Judge}, model memproses isi kedua jawaban sekaligus urutan penempatannya dan kalimat instruksi \textit{prompt} dalam satu distribusi probabilitas kondisional \cite{guSurveyLLMasaJudge2025}\cite{liLLMsasJudgesComprehensiveSurvey2024}. Oleh karena itu, pengujian yang hanya melaporkan akurasi mentah pada satu urutan jawaban dan satu \textit{prompt} belum cukup untuk menilai reliabilitas model. Nilai kecocokan mentah dapat meningkat akibat peluang kebetulan dan ketimpangan proporsi label \cite{raoAgreementMetricsLLMasJudge2026} sehingga memicu \textit{kappa deflation} \cite{normanReliabilityValiditySystematic2026}, sedangkan keputusan model tetap rentan berubah akibat \textit{position bias} \cite{shiJudgingJudgesSystematic}\cite{anghelDiagnosingBiasInstability2025} dan \textit{prompt sensitivity} \cite{weiSystematicEvaluationLLMasaJudge2025}\cite{bellibatluJudgeSenseBenchmarkPrompt2026}, terutama pada evaluasi lintas bahasa dan Bahasa Indonesia yang peka terhadap perubahan bahasa penyajian maupun bahasa instruksi \textit{prompt} \cite{fuHowReliableMultilingual2025}\cite{wiyonoIndoPrefMultiDomainPairwise2025}\cite{yinDoesJudgePrefer2026}.

Mengingat penelitian khusus Bahasa Indonesia berbasis \textit{human preference} yang menguji pengaruh posisi jawaban dan variasi \textit{prompt} secara terpisah masih terbatas, landasan teori tersebut digunakan sebagai dasar untuk menjawab rumusan masalah dan menguji hipotesis penelitian. Pertama, pengukuran kesepakatan yang memperhitungkan peluang kebetulan digunakan untuk menilai tingkat konsistensi \texttt{Qwen/Qwen3-32B} terhadap \textit{reference human preference} IndoPref melalui metrik \textit{accuracy} dan Cohen's kappa. Kedua, pengujian pertukaran posisi jawaban digunakan untuk menguji hipotesis pertama mengenai pengaruh perubahan urutan A--B dan B--A melalui \textit{position consistency}, \textit{preference flip rate}, dan \textit{first-position preference rate}. Ketiga, penggunaan tiga variasi \textit{prompt} yang bermakna sama (\textit{semantically equivalent}) digunakan untuk menguji hipotesis kedua mengenai pengaruh variasi \textit{prompt} P1, P2, dan P3 melalui \textit{prompt consistency} serta \textit{prompt flip rate} berpasangan, yang tahapannya dijelaskan lebih lanjut pada Bab~III.

\newpage
